%\clearpage

\section{Summary of changes included in version 10 of T2K-TN-141 (released on Sept. 18, 2013)}

This section is designed to allow readers to keep track of the changes between versions of T2K-TN-141.
Table, figure, equation, footnote or citation numbers shown in black refer to version of 9,
{\color{red}whereas the numbers shown in red refer to version of 10}.
Both numbers are shown where possible.

\begin{itemize}

\item
  Studies of fits using the 2-flavour approximation for the oscillation probabilities
  have been added as Appendix {\color{red}J}.
  All additional fits have been done with the
  '2+N' fitter, and use the Run 1+2+3 dataset only.
  \begin{itemize}
  \item The best-fit point has been added to Appendix {\color{red}J}.
  \item The pulls of the systematics at the best-fit point is shown in Fig.\ {\color{red}469}.
  \item The best-fit spectrum is shown in Fig.\ {\color{red}470} for the fine binning scheme, and
    Fig.\ {\color{red}471} for the coarse binning scheme.
  \item The $\Delta \chi^2$ surface is shown in Fig.\ {\color{red}472}.
  \item The contour made using the constant-$\Delta \chi^2$ method is shown in Fig.\ {\color{red}473}.
  \item A comparison between the 2-flavour contour with the two $\sin^22\theta_{23}$ 3-flavour fits
    is shown in Fig.\ {\color{red}474}.
  \end{itemize}

\item
  Studies of fits using $\sin^2\theta_{23}$ instead of $\sin^22\theta_{23}$ and 3-flavour oscillation 
  probabilities in vacuum instead of matter have been added as Appendix {\color{red}K}.
  All additional fits have been done with the
  '2+N' fitter, and use the Run 1+2+3 dataset only.
  \begin{itemize}
  \item The best-fit point has been added to Appendix {\color{red}K}.
  \item The pulls of the systematics at the best-fit point is shown in Fig.\ {\color{red}475}.
  \item The best-fit spectrum is shown in Fig.\ {\color{red}476} for the fine binning scheme, and
    Fig.\ {\color{red}477} for the coarse binning scheme.
  \item The $\Delta \chi^2$ surface is shown in Fig.\ {\color{red}478}.
  \item The contour made using the constant $\Delta \chi^2$ method is shown in Fig.\ {\color{red}479}.
  \item A comparison between the $\sin^2\theta_{23}$ contours with and without matter effects
    is shown in Fig.\ {\color{red}480}.
  \end{itemize}

\item Filled in the missing value of the spectra difference on line 877 {\color{red}(875)}.
\item Fixed the x-axis titles of Figs.\ 71 \& 88 {\color{red}(71 \& 88)}.

\end{itemize}

